% !TeX root = ../Book1.tex
\begin{figure}[htbp]
\centering
\begin{tikzpicture}[scale=1.02]
  \coordinate (v) at (0,3);
  \coordinate (w1) at (-2.3,0);
  \coordinate (w2) at (-0.6,-1);
  \coordinate (w3) at (2,-0.35);
  \coordinate (w4) at (1.5,0.8);
  \coordinate (p1) at (-1.035,1.65);
  \coordinate (p2) at (-0.27,1.20);
  \coordinate (p3) at (0.90,1.49);
  \coordinate (p4) at (0.675,2.01);
  \fill[color=AlgebraGray,opacity=0.13]
    (0.29,3.70)--(-2.585,-0.05)--(-2.315,-1.95)--(-0.13,0.90)--cycle;
  \draw[dashed,color=AlgebraGray]
    (0.29,3.70)--(-2.585,-0.05)--(-2.315,-1.95)--(-0.13,0.90)--cycle;
  \fill[color=AlgebraBlue,opacity=0.09]
    (-1.7,2.08)--(-1.27,0.92)--(1.35,1.12)--(1.43,2.36)--cycle;
  \draw[dashed,color=AlgebraBlue]
    (-1.7,2.08)--(-1.27,0.92)--(1.35,1.12)--(1.43,2.36)--cycle;
  \draw[gray] (w1)--(w2)--(w3)--(w4)--cycle;
  \draw[gray] (v)--(w1) (v)--(w2) (v)--(w3);
  \draw[dashed,gray] (v)--(w4);
  \draw[dotted,gray] (w1)--(-2.76,-0.60)
    (w2)--(-0.72,-1.80) (w3)--(2.40,-1.02)
    (w4)--(1.80,0.36);
  \draw[very thick,color=AlgebraBlue] (p1)--(p2)--(p3)--(p4)--cycle;
  \draw[dashed,line width=0.9pt]
    (-1.5705,1.965)--(0.30375,0.8625);
  \draw[line width=1.35pt] (p1)--(p2);
  \draw[->,line width=0.8pt] (-0.65,1.425)--(0.15,1.72);
  \foreach \p in {p1,p2,p3,p4} {\fill[color=AlgebraBlue] (\p) circle (1.4pt);}
  \foreach \p in {v,w1,w2,w3,w4} {\fill (\p) circle (1.4pt);}
  \node[above] at (v) {\(v\)};
  \node[left] at (w1) {\(w_1\)};
  \node[below] at (w2) {\(w_2\)};
  \node[right] at (w3) {\(w_3\)};
  \node[right] at (w4) {\(w_4\)};
  \node[left,color=AlgebraBlue] at (p1) {\(p_1\)};
  \node[below,color=AlgebraBlue] at (p2) {\(p_2\)};
  \node[below,color=AlgebraBlue] at (p3) {\(p_3\)};
  \node[right,color=AlgebraBlue] at (p4) {\(p_4\)};
  \node[color=AlgebraBlue] at (2.02,2.23) {\(\Pi\)};
  \node[anchor=south east] at (-1.55,2.06) {\(a(x-v)=r\)};
  \draw[thin,color=AlgebraGray] (0.17,1.72)--(1.92,1.65);
  \node[anchor=west] at (2.05,1.65) {\(g(x-v)<0\)};
  \node[color=AlgebraGray] at (-1.55,-1.65) {\(g(x-v)=0\)};
  \node[below] at (2.15,-1.20) {\(C+v\)};
\end{tikzpicture}
\caption[顶点截面的局部对应]{四棱锥的顶点截面及支持面的提升。蓝色截面位于 \(\Pi\)，灰色平面经过锥顶 \(v\) 与截面边 \([p_1,p_2]\)，其方程为 \(g(x-v)=0\)。截面内的黑色虚线为支持直线 \(a(x-v)=r\)，箭头指向 \(g(x-v)<0\) 的内侧；延长邻边所得的锥 \(C+v\) 全在相应闭半空间内。原面 \(F_{12}=\operatorname{conv}\{v,w_1,w_2\}\) 是灰色平面与原多面体的交。}
\label{b1:fig:vertex-section-local}
\end{figure}
